% 《大学物理简明教程》第5章气体动理论，解答题5.3.10—5.3.18。
% 用户扫描PDF/书内页126，保留模型、数据与编号，独立LaTeX推导；待学术二校。
\subsection{解答题：能量、分布函数与碰撞自由程}

\sourceproblem{5.3.10}{126}{126}
\problemtext 一定量刚性双原子分子理想气体体积为
$V=2\times10^{-3}\,\mathrm{m^3}$，内能$U=6.75\times10^2$ J。
(1) 求压强；(2) 若总分子数$N=5.4\times10^{22}$，
求每分子平均平动动能和温度。
\workings
刚性双原子分子自由度$i=5$，由
$U=\tfrac52Nk_{\mathrm B}T$及$pV=Nk_{\mathrm B}T$，
\[
\boxed{p=\frac{2U}{5V}
=\frac{2(675)}{5(0.002)}
=1.35\times10^5\,\mathrm{Pa}}.
\]
平均平动动能
\[
\boxed{\overline\varepsilon_{\mathrm{trans}}
=\frac32k_{\mathrm B}T=\frac{3U}{5N}
=7.50\times10^{-21}\,\mathrm J}.
\]
温度
\[
\boxed{T=\frac{2U}{5Nk_{\mathrm B}}
\approx362\,\mathrm K}.
\]
\crosscheck $pV=(2/5)U$，两边均有焦耳量纲；
计算温度不需要知道气体分子质量。

\sourceproblem{5.3.11}{126}{126}
\problemtext 容积$1\,\mathrm{m^3}$的容器内有1 mol氢气（按刚性双原子分子）；
整个容器原以$10\,\mathrm{m/s}$平移，突然停下。
若氢气整体运动的机械动能中有80\%转化成分子热运动内能，求气体温度与压强的升高量。
\workings
1 mol氢气质量$m_{\mathrm H_2}\approx0.002\,\mathrm{kg}$，
整体平动动能
$\tfrac12mv^2=\tfrac12(0.002)(10^2)=0.10$ J。
进入气体内能的部分为
$\Delta U=0.8(0.10)=0.08$ J。
由$n=1$ mol与定容热容$C_V=\tfrac52R$，
\[
\boxed{\Delta T=\frac{\Delta U}{(5/2)R}
\approx3.85\times10^{-3}\,\mathrm K}.
\]
再利用$pV=nRT$，容积固定，
\[
\boxed{\Delta p=\frac{nR\Delta T}{V}
=\frac{2}{5}\frac{\Delta U}{V}
=0.032\,\mathrm{Pa}}.
\]
\crosscheck 如把80\%的转化误当100\%，压强会被高估至0.040 Pa；
气体升温虽很小，但不应为零。

\sourceproblem{5.3.12}{126}{126}
\problemtext 标准状态下氧气（刚性双原子分子）和氦气的体积之比
$V_{\mathrm{O_2}}/V_{\mathrm{He}}=1/2$，求两种气体内能之比。
\workings
标准状态两气体具有相同$T,p$，故$n_{\mathrm{O_2}}/n_{\mathrm{He}}
=V_{\mathrm{O_2}}/V_{\mathrm{He}}=1/2$。
氧气$i=5$，氦气$i=3$，
\[
\boxed{\frac{U_{\mathrm{O_2}}}{U_{\mathrm{He}}}
=\frac{(5/2)n_{\mathrm{O_2}}RT}
 {(3/2)n_{\mathrm{He}}RT}
=\frac{5}{6}}.
\]
\crosscheck 总内能比较既要计入分子自由度，又要计入不同容积导致的摩尔数差异。

\sourceproblem{5.3.13}{126}{126}
\problemtext 说明麦克斯韦速率分布函数$f(v)$的意义，并解释
$f(v)\mathrm dv$、$nf(v)\mathrm dv$、$Nf(v)\mathrm dv$、
$\int_0^v f(s)\mathrm ds$、$\int_0^\infty f(s)\mathrm ds$及
$\int_{v_1}^{v_2}Nf(v)\mathrm dv$。
其中$n$为分子数密度，$N$为整个系统的分子总数。
\workings
$f(v)$为每单位速率的\emph{概率密度}，单位为$\mathrm{s/m}$。
\[
\begin{array}{l|l}
\text{表达式}&\text{物理意义}\\\hline
f(v)\mathrm dv&\text{速率落在}[v,v+\mathrm dv]\text{的分子所占比例}\\
nf(v)\mathrm dv&\text{单位体积内该速率区间的分子数}\\
Nf(v)\mathrm dv&\text{整个系统内该速率区间的分子数}\\
\int_0^v f(s)\mathrm ds&\text{速率不大于}v\text{的分子比例}\\
\int_0^\infty f(s)\mathrm ds&1\text{（归一化）}\\
\int_{v_1}^{v_2}Nf(v)\mathrm dv&
\text{系统内速率在}[v_1,v_2]\text{的分子数}
\end{array}
\]
\crosscheck 只有概率、比例无量纲；分子数在统计模型中可按连续分布近似为期望数量。

\sourceproblem{5.3.14}{126}{126}
\problemtext 质量$m=6.2\times10^{-17}\,\mathrm{kg}$的微粒悬浮于
$27^\circ\mathrm C$的液体，观测到其方均根速率为$1.4\,\mathrm{cm/s}$。
若按经典麦克斯韦分布求伏尔伽德罗常量（题面印作伏尔伽德罗常量，实际应为玻尔兹曼常量），应得多少？
\workings
核对原题末句为“求阿伏伽德罗常量”，而仅由$m,v_{\mathrm{rms}},T$
能得到的是\emph{玻尔兹曼常量}$k_{\mathrm B}$；
若要换算阿伏伽德罗常量，必须另给气体常量$R$。
令$v_{\mathrm{rms}}=0.014\,\mathrm{m/s}$、$T=300.15\,\mathrm K$，
由$\tfrac12m\langle v^2\rangle=\tfrac32k_{\mathrm B}T$得
\[
\boxed{k_{\mathrm B}=\frac{m v_{\mathrm{rms}}^2}{3T}
=\frac{(6.2\times10^{-17})(0.014)^2}{3(300.15)}
\approx1.35\times10^{-23}\,\mathrm{J/K}}.
\]
取$R=8.314462618\,\mathrm{J/(mol\,K)}$，
则
\[
\boxed{N_A=\frac{R}{k_{\mathrm B}}
\approx6.16\times10^{23}\,\mathrm{mol^{-1}}}.
\]
\crosscheck 本实验估计值与标准$6.022\times10^{23}\,\mathrm{mol^{-1}}$同数量级；
估算偏差来自给定实验数值舍入。原题需再核对其所称常量的中文用字。

\sourceproblem{5.3.15}{126}{126}
\problemtext 两瓶氧气、氢气处于相同压强与绝对温度，
氧气体积是氢气体积的两倍。
求两种气体的分子数密度之比和方均根速率之比。
\workings
理想气体分子数密度$n=p/(k_{\mathrm B}T)$，
故尽管两容器体积不同，
\[
\boxed{\frac{n_{\mathrm{O_2}}}{n_{\mathrm{H_2}}}=1}.
\]
对于方均根速率，
$v_{\mathrm{rms}}=\sqrt{3RT/M_{\mathrm{mol}}}$，故
\[
\boxed{\frac{v_{\mathrm{rms,O_2}}}{v_{\mathrm{rms,H_2}}}
=\sqrt{\frac{M_{\mathrm{H_2}}}{M_{\mathrm{O_2}}}}
=\sqrt{\frac{2}{32}}=\frac14}.
\]
\crosscheck 容积比2:1使\emph{总分子数}比为2:1，但\emph{数密度}仍为1:1，不得混淆。

\sourceproblem{5.3.16}{126}{126}
\problemtext $N$个粒子的速率分布图中，$Nf(v)$由原点线性升至
$(v_0,a)$，然后在$v_0\le v\le2v_0$恒等于$a$，区间外为零。
求$f(v)$、平均速率与速率位于$1.5v_0$到$2v_0$间的粒子数。
\workings
归一化要求图下面积为$N$：
\[
N=\frac12av_0+av_0=\frac32av_0
\Longrightarrow a=\frac{2N}{3v_0}.
\]
故
\[
\boxed{f(v)=
\begin{cases}
\frac{2v}{3v_0^2},&0\le v<v_0,\\
\frac{2}{3v_0},&v_0\le v\le2v_0,\\
0,&\text{其他}.
\end{cases}}
\]
平均速率
\[
\overline v=\int_0^\infty v f(v)\,dv
=\frac{2}{3v_0^2}\frac{v_0^3}{3}
+\frac{2}{3v_0}\frac{(2v_0)^2-v_0^2}{2}
=\boxed{\frac{11}{9}v_0}.
\]
所求区间的分子数
\[
\boxed{N_{[1.5v_0,2v_0]}
=\int_{1.5v_0}^{2v_0}Nf(v)\,dv
=\frac N3}.
\]
\crosscheck $\int f(v)\,dv=1$；区间概率为$1/3$不随$v_0$选取而变化。

\sourceproblem{5.3.17}{126}{126}
\problemtext 真空管内压强$p=1.38\times10^{-3}\,\mathrm{Pa}$，
温度$27^\circ\mathrm C$；气体分子有效直径$d=3.0\times10^{-10}\,\mathrm m$。
求分子数密度和平均自由程。
\workings
温度$T=300.15$ K，
\[
\boxed{n=\frac p{k_{\mathrm B}T}
\approx3.33\times10^{17}\,\mathrm{m^{-3}}}.
\]
硬球气体自由程
\[
\boxed{\overline\lambda
=\frac{1}{\sqrt2\pi d^2n}
=\frac{k_{\mathrm B}T}{\sqrt2\pi d^2p}
\approx7.51\,\mathrm m}.
\]
\crosscheck 当管道尺寸远小于计算自由程时，
实际粒子与器壁的碰撞可能远比粒子间碰撞频繁；
这里求的是硬球\emph{分子间}平均自由程。

\sourceproblem{5.3.18}{126}{126}
\problemtext 标准状态下氢气分子有效直径
$d=2.30\times10^{-10}\,\mathrm m$。
求平均自由程和平均碰撞频率。
\workings
取$p_0=101325$ Pa、$T_0=273.15$ K、
$M_{\mathrm{H_2}}=2.0\times10^{-3}\,\mathrm{kg/mol}$。
数密度
$n_0=p_0/(k_{\mathrm B}T_0)\approx2.687\times10^{25}\,\mathrm{m^{-3}}$。
\[
\boxed{\overline\lambda=
\frac{1}{\sqrt2\pi d^2n_0}
\approx1.58\times10^{-7}\,\mathrm m}.
\]
按经典麦克斯韦分布
$\overline v=\sqrt{8RT_0/(\pi M_{\mathrm{H_2}})}
\approx1.70\times10^3\,\mathrm{m/s}$，故
\[
\boxed{\overline Z=\frac{\overline v}{\overline\lambda}
\approx1.07\times10^{10}\,\mathrm{s^{-1}}}.
\]
\crosscheck $\overline v$必须用\emph{平均速率}而非方均根速率计算平均碰撞频率；不同的氢气摩尔质量舍入会引入约1\%量级的末位差异。
